\subsection{Examples}

\begin{enumerate}
\def\labelenumi{\arabic{enumi})}
\item
  Let E be a normed space. Equip E with the product defined by ab = 0 for all a and b in E; this defines on E a normed algebra structure.
\item
  Let X be a set. We denote by \(\mathcal{B}(X;K)\) the normed algebra of bounded functions on X with values in K, equipped with the norm

\[
\| f \| = \sup _ {x \in X} | f (x) |
\]

(TG, X, p.~22). It is a unital Banach algebra over K (TG, X, p.~21, cor. 1). It is commutative. Let \(f\) be an element of \(\mathcal{B}(\mathrm{X};\mathrm{K})\) . Then \(f\) is invertible in \(\mathcal{B}(\mathrm{X};\mathrm{K})\) if and only if one has

\[
\inf _ {x \in X} | f (x) | > 0.
\]

The spectrum of \(f\) is therefore the set of \(\lambda\in K\) such that

\[
\inf _ {x \in \mathrm{X}} | f (x) - \lambda | = 0,
\]

that is, the closure in K of the set \(f(\mathrm{X})\) of values of \(f\) .

\item
  Let X be a topological space. We denote by \(\mathcal{C}_{b}(\mathrm{X};\mathrm{K})\) the unital subalgebra of \(\mathcal{B}(\mathrm{X};\mathrm{K})\) of continuous and bounded functions on X with values in K, \(\mathcal{C}_{0}(\mathrm{X};\mathrm{K})\) the subalgebra of \(\mathcal{C}_{b}(\mathrm{X};\mathrm{K})\) consisting of functions that tend to 0 at infinity (cf.~INT, III, §1, n°2 and TG, X, p.~40, cor. 2). We recall that \(\mathcal{H}(\mathrm{X};\mathrm{K})\) denotes the subalgebra of \(\mathcal{C}_{b}(\mathrm{X};\mathrm{K})\) of continuous functions with compact support.
The algebras \(\mathcal{C}_{b}(\mathrm{X};\mathrm{K})\) and \(\mathcal{C}_{0}(\mathrm{X};\mathrm{K})\) are commutative Banach algebras over K; indeed, they are closed subspaces of \(\mathcal{B}(\mathrm{X};\mathrm{K})\) (TG, X, p.~21, cor. 2 and INT, III, §1, n°2).

Since the inverse of a continuous function that is nowhere zero is continuous, the subalgebra \(\mathcal{C}_{b}(\mathrm{X};\mathrm{K})\) is a full subalgebra of \(\mathcal{B}(\mathrm{X};\mathrm{K})\) . In particular, the spectrum of an element f of \(\mathcal{C}_{b}(\mathrm{X};\mathrm{K})\) is equal to \(\overline{f(\mathrm{X})}\) .

If X is discrete, one has \(\mathcal{C}_b(\mathrm{X};\mathrm{K}) = \mathcal{B}(\mathrm{X};\mathrm{K})\) .

If X is compact, one has \(\mathcal{C}_{b}(\mathrm{X};\mathrm{K}) = \mathcal{H}(\mathrm{X};\mathrm{K}) = \mathcal{C}(\mathrm{X};\mathrm{K})\) , the unital normed algebra \(\mathcal{C}(\mathrm{X};\mathrm{K})\) of continuous functions on X with values in K. In this case, an element \(f \in \mathcal{C}(\mathrm{X};\mathrm{K})\) is invertible if and only if f does not take the value 0, and the spectrum of f is equal to the set \(f(\mathrm{X})\) of values of f.

Suppose henceforth that X is not compact. The algebra \(\mathcal{C}_{0}(X;K)\) is then not unital; the algebra obtained from it by adjoining a unit element is identified with the subalgebra \(K \cdot 1 \oplus \mathcal{C}_{0}(X;K)\) of \(\mathcal{C}(X;K)\) formed by functions that have a limit at infinity. For an element of this subalgebra to be invertible, it is necessary and sufficient that it does not vanish and that its limit at infinity is not zero. It follows that for every \(f \in \mathcal{C}_{0}(X;K)\) , the spectrum of f is equal to \(f(X) \cup \{0\}\) .

The spectrum of \(f \in \mathcal{K}(\mathrm{X};\mathrm{K})\) is equal to \(f(\mathrm{X})\) (indeed, if \(\mathrm{X}\) is not compact, 0 belongs to \(f(\mathrm{X})\) ).

\item
  Let \(n \geqslant 0\) be an integer. Let \(A_{n}\) be the algebra of functions \(f: [0,1] \to K\) admitting continuous derivatives in \([0,1]\) up to order \(n\) , equipped with the norm
\[
\| f \| = \sum_ {k = 0} ^ {n} \frac {1}{k !} \sup _ {0 \leqslant t \leqslant 1} | f ^ {(k)} (t) |.
\]

If \(f,g\in \mathrm{A}_n\) , one has

\[
\begin{array}{l} \| f g \| = \sum_ {k = 0} ^ {n} \frac {1}{k !} \sup | (f g) ^ {(k)} (t) | = \sum_ {k = 0} ^ {n} \frac {1}{k !} \sup \Bigl | \sum_ {s = 0} ^ {k} \binom {k} {s} f ^ {(s)} (t) g ^ {(k - s)} (t) \Bigr | \\ \leqslant \sum_ {k = 0} ^ {n} \sum_ {s = 0} ^ {k} \frac {1}{s ! (k - s) !} \sup _ {0 \leqslant t \leqslant 1} | f ^ {(s)} (t) | \sup _ {0 \leqslant t \leqslant 1} | g ^ {(k - s)} (t) | = \| f \|   \| g \|, \end{array}
\]

by Leibniz's formula (FVR, I, p.~28, prop. 2), hence \(A_{n}\) is a unital commutative Banach algebra.

\item
  Let E be a Banach space whose norm we denote by p. The algebra \(\mathcal{L}(\mathrm{E})\) of continuous endomorphisms of E, equipped with the norm
\[
\| u \| = \sup _ {p (x) \leqslant 1} p (u (x))
\]

is a unital Banach algebra (EVT, III, p.~14 and p.~24, cor. 2; TG, X, p.~23, formula (3)).

\item
  Let G be a locally compact group. Denote by e its identity element. Let \(\mathcal{M}^{1}(G)\) be the Banach space of bounded complex measures on G (INT, III, p.~57). The convolution product (INT, VIII, p.~120, def. 1) equips \(\mathcal{M}^{1}(G)\) with a complex Banach algebra structure (INT, VIII, §3, n° 1, prop. 2) admitting as identity element the measure \(\varepsilon_{e}\) defined by the unit mass placed at the point e. If G is commutative, this Banach algebra is commutative. The space \(\mathcal{C}'(G)\) of measures with compact support is a subalgebra of \(\mathcal{M}^{1}(G)\) (loc. cit.).
\item
  Let G be a locally compact group equipped with a Haar measure \(\mu\) . Then \(\mathrm{L}_{\mathrm{K}}^{1}(\mathrm{G},\mu)\) is a Banach algebra for the convolution product (INT, VIII, prop. 12, p.~166). If K = C, the map defined by \(f \mapsto f \mu\) allows one to identify \(\mathrm{L}_{\mathbf{C}}^{1}(\mathrm{G},\mu)\) with a subalgebra of the Banach algebra \(\mathcal{M}^{1}(\mathrm{G})\) .
If G is commutative, the Banach algebra \(\mathrm{L}_{\mathrm{K}}^{1}(\mathrm{G},\mu)\) is commutative.

\item
  Take G = Z and K = C in Example 7. Then \(\mathrm{L}_{\mathbf{C}}^{1}(\mathbf{Z})\) is the complex commutative Banach algebra of sequences \((x_{n})_{n\in\mathbf{Z}}\) such that \(\sum_{n}|x_{n}|<+\infty\) , the product of elements \((x_{n})\) and \((y_{n})\) being \((z_{n})\) , where
\[
z _ {n} = \sum_ {k \in {\bf Z}} x _ {k} y _ {n - k}
\]

and the norm \(\|(x_{n})\|=\sum_{n}|x_{n}|\) . This algebra admits as its identity element the sequence \(\varepsilon=(\varepsilon_{n})\) such that \(\varepsilon_{0}=1\) and \(\varepsilon_{n}=0\) for \(n\neq0\) .

Let U denote the unit circle in C. If \(x = (c_{n})\) is an element of A, let \(\varphi(x)\) be the continuous function on U whose value at \(e^{it}\) is

\[
\varphi (x) (e ^ {i t}) = \sum_ {n \in {\bf Z}} c _ {n} e ^ {i n t}.
\]

One checks that \(\varphi\) is a morphism of \(\mathrm{L}_{\mathbf{C}}^{1}(\mathbf{Z})\) onto an algebra A of continuous functions on \(\mathbf{U}\) , multiplication in A being the usual multiplication. Integrating the equality term by term

\[
\left(\sum_ {m \in {\bf Z}} c _ {m} e ^ {i m t}\right) \cdot e ^ {- i n t} = \varphi (x) (e ^ {i t}) \cdot e ^ {- i n t},
\]

it comes

\[
c _ {n} = \frac {1}{2 \pi} \int_ {0} ^ {1} \varphi (x) (e ^ {i t}) e ^ {- i n t} d t.
\]

In particular, it follows that the morphism \(\varphi\) is injective. The algebra A, endowed with the norm induced by that of \(\mathrm{L}_{\mathbf{C}}^{1}(\mathbf{Z})\) via \(\varphi\) , is called the Banach algebra of absolutely convergent Fourier series. It admits as identity element the function \(1 = \varphi(\varepsilon)\) .

\item
  Let \(\Delta\) be the disk of complex numbers \(z\) satisfying \(|z| \leqslant 1\) . The algebra A of continuous functions on \(\Delta\) that are analytic in the interior of \(\Delta\) (VAR, R1, p.~26, 3.2.1) is equipped with the norm \(\|f\| = \sup_{z \in \Delta} |f(z)|\) . Then A is a commutative unital Banach algebra.
\end{enumerate}

